\begin{statement}[C.7.9; Proved: the endomorphism and slope obstructions]
Let $E/\overline\Q$ be an elliptic curve, let
$q\in E^\vee(\overline\Q)$ be nontorsion, and let
\[
 1\longrightarrow\mathbb G_m\longrightarrow G_q
       \longrightarrow E\longrightarrow0
\]
be the extension of class $q$. Then
$\operatorname{End}_{\overline\Q}(G_q)=\Z$, including when
$E$ has complex multiplication. In particular a noninteger
CM endomorphism of $E$ does not lift to an endomorphism of this
fixed extension.

For
\[
 E:y^2+y=x^3-38x+90,\qquad P=(0,9),\qquad Q=2P=(4,-2),
\]
one has
\[
 N(Q)=19^{1/4}N(P)^4.
\]
Use the differential $dx/(2y+1)$, choose the positive real
period $\omega_1$, and choose the real logarithm
$u\in(0,\omega_1/2)$ of $Q$. Put
\begin{equation}\label{eq:cd-slope}
 c=\zeta(u)-\frac{\eta_1}{\omega_1}u,
 \qquad \eta_1=\eta(\omega_1).
\end{equation}
Then $c$ is a real transcendental number. If $N(P)$ were
algebraic, the extension $G_Q$ of C.7.1 would have a
nontorsion algebraic point whose logarithm is $(u,cu)$.
The complex line
$\ell_c=\{(z,cz):z\in\C\}$ is not defined over
$\overline\Q$ in the algebraic logarithm coordinates of
C.7.1. Thus the analytic subgroup theorem cannot be applied
to that line using the hypothesized compact point.
\end{statement}
\begin{proof}
The torus is the maximal connected affine subgroup of $G_q$.
Consequently every endomorphism of $G_q$ induces multiplication
by some $n\in\Z$ on $\mathbb G_m$ and an endomorphism
$\alpha$ on $E$. Functoriality of extension classes gives
\begin{equation}\label{eq:cd-compatibility}
 \alpha^*q=nq\quad\text{in }E^\vee.
\end{equation}
Here $\alpha^*$ is the dual endomorphism, not $\alpha$
acting on an arbitrarily identified copy of $E$. To make the
convention explicit, choose a principal polarization
$\lambda:E\to E^\vee$, put $Q_0=\lambda^{-1}q$, and
write $\alpha^\dagger=\lambda^{-1}\alpha^*\lambda$.
Equation \eqref{eq:cd-compatibility} says
$(\alpha^\dagger-[n])Q_0=O$. A nonzero endomorphism of an
elliptic curve is an isogeny and has finite kernel. Since
$Q_0$ is nontorsion, $\alpha^\dagger=[n]$, hence
$\alpha=[n]$ as well.

The difference between the given endomorphism and $[n]$ has
image in the torus and is trivial on the torus. It therefore
factors through an algebraic-group homomorphism
$E\to\mathbb G_m$. This homomorphism is zero, since a
morphism from the proper connected curve $E$ to the affine
group $\mathbb G_m$ is constant and a group homomorphism
preserves the identity. Thus there is no additional ambiguity
from $\operatorname{Hom}(E,\mathbb G_m)$, and the
endomorphism is exactly $[n]$. Conversely all integer
multiplications exist. This proves the first assertion.

For the concrete curve, the tangent slope at $P$ is $-2$,
giving $2P=(4,-2)$. Its discriminant is $-19^3$ and
$\psi_2(P)=2\cdot9+1=19$. The division identity in C.4
therefore gives the asserted formula for $N(Q)$, so replacing
$P$ by $Q$ preserves algebraicity of the norm. This doubling
also places the base point in the identity component of the
real elliptic group without any component assumption. In the
chosen differential the Weierstrass coordinates are
$(\wp,\wp')=(x,2y+1)$, with $g_2=152$, $g_3=-361$.
The negative discriminant means that the real cubic has one
real root. Along $(0,\omega_1/2)$, $\wp$ decreases from
$+\infty$ to that root and $\wp'<0$. Since
$\wp'(u)=2(-2)+1=-3$, the indicated real logarithm of $Q$
exists.

The period vectors of $G_Q$ in C.7.1 include
\[
 (\omega_1,\zeta(u)\omega_1-\eta_1u)
       =\omega_1(1,c).
\]
These are algebraic logarithm coordinates: $dz$ and the
invariant differential $d\log r-df_u/f_u$ are algebraic,
because $df_u/f_u$ is the algebraic third-kind differential
displayed in C.7.1. Real invariants and real $u,\omega_1$
make $c$ real.

Suppose $c$ were algebraic. Then $\ell_c$ would be an
algebraic Lie subspace. Dividing the nonzero period
$\omega_1(1,c)$ by a sufficiently large positive integer
gives a vector in $\ell_c$ whose exponential is a nonzero
torsion point of $G_Q$, hence an algebraic point. The
analytic subgroup theorem
\citep[Theorem 6.1]{huber2022} supplies a positive-dimensional
connected algebraic subgroup with Lie algebra contained in
$\ell_c$. Its dimension is one, and its projection to $E$
is nontrivial. Such a subgroup must be an elliptic curve
$E'$ projecting isogenously to $E$: a connected
one-dimensional affine algebraic group cannot map
nontrivially to an elliptic curve. Its inclusion splits the
pullback of $G_Q$ to $E'$. The dual of that isogeny would
therefore kill the extension class $\lambda(Q)$, forcing
$Q$ torsion. This contradicts C.6.R1. We conclude that $c$
is transcendental. In particular the line is not defined over
$\overline\Q$. This argument uses the period vector alone;
it does not assume algebraicity of $N(P)$.

Finally assume, only for the attempted contradiction, that
$N(P)$ is algebraic. The positive real number
\[
 r=3\,19^{-1/2}N(Q)^2
\]
is algebraic, and C.7.2 places the point $R=(Q,r)$ of
$G_Q$ in its maximal compact subgroup. The chosen real
logarithm has
\[
 f_u(u)=-\wp'(u)\sigma(u)^2e^{-u\zeta(u)}
        =3\sigma(u)^2e^{-u\zeta(u)}>0.
\]
Consequently $v=\log(r/f_u(u))$ is real and
$(u,v)$ is a logarithm of $R$, with no imaginary logarithm
ambiguity. Substituting this real logarithm into the exact
compactness equation $\chi=0$ in C.7.2 gives
\[
 v=\zeta(u)u-u\eta(u)
   =u\left(\zeta(u)-\frac{\eta_1}{\omega_1}u\right)=cu.
\]
Its projection $Q$ is nontorsion, so $R$ is nontorsion.
The input missing from this proposed use of the analytic
subgroup theorem is a proper Lie subspace defined over
$\overline\Q$; the actual compact slope line is proved
not to have that property. The CM decomposition makes the
same obstruction visible as
\[
 c=\zeta(u)-s_2u-\frac{\pi}{A}u.
\]
The last term is the derivative of the real Gaussian
correction. It cannot be discarded while retaining the
compactness equation.
\end{proof}
